\documentclass[11pt]{article}
\usepackage[margin=1in]{geometry}
\usepackage{amsmath,amssymb,amsthm}
\usepackage[hidelinks]{hyperref}
\newtheorem{proposition}{Proposition}
\newtheorem{lemma}[proposition]{Lemma}
\newcommand{\Lfac}{\mathcal L}
\newcommand{\C}{\mathbb C}
\newcommand{\A}{\mathcal A}
\begin{document}
\title{Unrestricted Strong Factorial: two mechanism obstructions}
\author{Internal route audit; no priority or breakthrough claim}
\date{18 November 2026}
\maketitle

\noindent\textbf{Target and significance.}
For every $d\geq1$, every nonzero $F=\sum_{j=1}^s c_jx^{a_j}$
with distinct $a_j\in\mathbb N^d$ and $c_j\ne0$, and every $n\geq1$,
unrestricted SFC asks that
\[
  \bigl(\Lfac(F^n),\ldots,\Lfac(F^{n+s-1})\bigr)\ne0,
  \qquad \Lfac(x^\alpha)=\alpha!.
\]
This finite-window assertion is stronger than FC or absence of an
infinite zero tail. It remains open here. A general proof or an actual
counterexample would meet the research objective; the route results
below do not.

The primary definition and rigidity bridge are Edo--van den Essen,
\href{https://arxiv.org/abs/1304.3956v2}{Definitions 2.7--2.8 and
Theorem 2.25}. Rocks, \emph{J. Algebra} 478 (2017), 341--352,
\href{https://people.math.sc.edu/rocks/BradyRocksSFC2015.pdf}{author
preprint, Theorems 3.3 and 3.10}, proves powers of linear forms and
powers of sums of same-prime-power variable monomials; Section 3.4
gives partial two-monomial results. These do not state an arbitrary
support-size window theorem.

The repository's degree-six and selected degree-seven binomial results
retain their exact scopes. Full Furter rigidity is already available
as a credited internal reproduction in
\href{RIGIDITY_CRITICAL_VALUES.tex}{the critical-value audit}, using
liqsweep's pinned September 2026 argument. Its factorial transfer concerns
$(\prod_i y_i)\sum_i\mu_i y_i$. Reproving rigidity or extending one
degree-seven terminal family would not resolve the present target.

\section{Monomial count does not bound scalar recurrence order}
A natural proposed proof would produce a homogeneous recurrence of
order at most $s$, with a nonzero backward coefficient, and propagate
an $s$-term zero block. Even the first requirement fails.

\begin{proposition}\label{prop:recurrence}
For the ordinary two-monomial phase $F=1+X^2$, put
\[
 M_n=\Lfac(F^n)=\sum_{k=0}^n\binom nk(2k)!.
\]
The minimal homogeneous recurrence order over $\C(n)$ is three,
including recurrences required only for all sufficiently large integers.
\end{proposition}
\begin{proof}
Integration by parts, with the boundary term retained, gives
\[
 M_n-2n(2n-1)M_{n-1}+4n(n-1)M_{n-2}=1,\qquad n\geq2.       \tag{1}
\]
Indeed $M_n=1+2nJ_n$, where
$J_n=\int_0^\infty x(1+x^2)^{n-1}e^{-x}\,dx$, and a second integration
by parts gives
$J_n=(2n-1)M_{n-1}-2(n-1)M_{n-2}$.
Subtracting (1) at $n-1$ gives the homogeneous order-three relation
\[
 M_n-(4n^2-2n+1)M_{n-1}
 +(n-1)(8n-6)M_{n-2}
 -4(n-1)(n-2)M_{n-3}=0.                                \tag{2}
\]
All following rational-function identities are understood away from
their finitely many poles.

The sequence cannot eventually be a rational function of $n$, since
$M_n\geq(2n)!$. If a homogeneous order-two relation existed, eliminating
one term against (1) would either contradict (1), make $M_n$ rational,
or give
\[
 M_n=r(n)M_{n-1}+s(n),\qquad r,s\in\C(n).
                                                               \tag{3}
\]
The same conclusion includes a homogeneous order-one relation.
Here $r\ne0$, again by growth. Put
$A(n)=2n(2n-1)$ and $B(n)=4n(n-1)$. Substitute (3) and its shift in (1).
The coefficient of $M_{n-2}$ is
\[
 (r(n)-A(n))r(n-1)+B(n).
\]
It must vanish identically, since otherwise $M_{n-2}$ is rational.
Thus $R(n)=r(n)/(2n)$ would satisfy
\[
 R(n)+\frac1{R(n-1)}=2n-1.                              \tag{4}
\]
Every finite zero of $R$ forces a pole of the same multiplicity one
step later, and every finite pole forces a zero one step earlier:
these poles must cancel in (4). Hence the reduced numerator and
denominator have equal degree. The left side then has a finite limit
at infinity, contradicting the linear right side.
Order zero is excluded by positivity. Equation (2) therefore has
minimal order three.
\end{proof}

This is not an SFC counterexample: every $M_n$ here is positive.
The inhomogeneous order-two relation (1) remains valid. The result
does not exclude coupled contiguity systems or nonlinear arguments;
it excludes the proposed homogeneous scalar order bound.

\section{All polynomial symmetries of the functional}
A second counterexample strategy would transform a known periodic
moment phase by a nonlinear polynomial symmetry of $\Lfac$, reducing
its number of monomials. In fact no such symmetry exists.

For $\A=\C[x_1,\ldots,x_d]$ let
\[
 \mathfrak g=\{D\in\operatorname{Der}_{\C}(\A):
                 \Lfac(Df)=0\text{ for every }f\in\A\}.
\]
This is a Lie algebra. For $P=\sum_\alpha p_\alpha x^\alpha$ put
$T(P)=\sum_\alpha p_\alpha\alpha!u^\alpha$.

\begin{lemma}\label{lem:derivations}
For $D=\sum_iP_i\partial_i$,
\[
 D\in\mathfrak g
 \quad\Longleftrightarrow\quad
 x_i\mid P_i\text{ for every }i,\qquad
 \sum_i(P_i-\partial_iP_i)=0.                            \tag{5}
\]
\end{lemma}
\begin{proof}
Formally
$Z(t)=\Lfac(e^{t\cdot x})=\prod_i(1-t_i)^{-1}$.
Set $u_i=(1-t_i)^{-1}$. Coefficientwise differentiation gives
$\Lfac(Pe^{t\cdot x})/Z(t)=T(P)(u)$.
Consequently $D\in\mathfrak g$ is equivalent to the Laurent identity
\[
 \sum_i(1-u_i^{-1})T(P_i)(u)=0.                          \tag{6}
\]
A negative exponent of $u_i$ can arise only from
$-u_i^{-1}T(P_i)|_{u_i=0}$. Its coefficient must be zero.
Since $T$ rescales every coefficient by a nonzero factorial,
this is exactly $x_i\mid P_i$. Under that divisibility,
$u_i^{-1}T(P_i)=T(\partial_iP_i)$.
The remaining identity is (5). These steps reverse, proving sufficiency.
Only formal power series and Laurent polynomials are used.
\end{proof}

\begin{lemma}\label{lem:boundary}
For $d\geq2$, define
$V(p)=\operatorname{span}_{\C}\{D(p):D\in\mathfrak g\}$.
Its rank-drop locus is precisely $\bigcup_i\{x_i=0\}$.
\end{lemma}
\begin{proof}
On $p_i=0$, every $D$ has zero $i$th component by (5), so the rank
drops. Conversely let every $p_i\ne0$. Given $i$, choose $j\ne i$ and put
\[
 h=\frac{x_j-p_j}{p_ip_j},\quad R=x_ix_jh,\quad
 P_i=\partial_jR-R,\quad P_j=R-\partial_iR,
\]
with other components zero.
Both divisibilities in (5) hold and
$\sum_k\partial_kP_k=\partial_jR-\partial_iR=\sum_kP_k$.
This gives $D\in\mathfrak g$. At $p$, $h=0$,
$\partial_jh=1/(p_ip_j)$ and $\partial_ih=0$, so $D(p)=e_i$.
Thus $V(p)=\C^d$ off the coordinate boundary.
\end{proof}

\begin{proposition}\label{prop:automorphisms}
For every $d\geq1$, the $\C$-algebra automorphisms $\phi$ of $\A$
satisfying $\Lfac(\phi(f))=\Lfac(f)$ for every polynomial $f$ are
exactly coordinate permutations.
\end{proposition}
\begin{proof}
Such an automorphism conjugates $\mathfrak g$ to itself:
$\Lfac(\phi D\phi^{-1}f)=\Lfac(D\phi^{-1}f)=0$.
Its Jacobian is invertible everywhere, so it preserves the rank-drop
locus in Lemma~\ref{lem:boundary}. It permutes the irreducible coordinate
hyperplanes. Therefore
$\phi(x_i)=c_ix_{\sigma(i)}$, with $c_i\ne0$ and $\sigma$ a permutation.
First moments give $c_i=1$. The converse is immediate.

For $d=1$, extend $\phi$ by the identity on one extra variable.
Product factorial moments make the extension preserve $\Lfac$ in two
variables. The preceding case forces the extension to be a permutation;
since it fixes the extra coordinate, $\phi$ is the identity.
\end{proof}

In particular, no nonzero locally nilpotent derivation preserves
$\Lfac$: its finite exponential would give a nonconstant algebraic
one-parameter subgroup of this finite permutation group.
These results concern preservation on \emph{every polynomial}.
Vanishing moments on $\C[F]$ or a finite zero block do not imply that
hypothesis, so the rigidity proposition does not itself prove FC or SFC.

\section{Equal valuations at the sharp window bound}
For a primitive $s$th root $\zeta$, the actual ordinary phase
\[
 F=\sum_{j=0}^{s-1}\zeta^jx_j
\]
has
\[
 \sum_{m\geq0}\frac{\Lfac(F^m)}{m!}t^m
 =\prod_{j=0}^{s-1}(1-\zeta^jt)^{-1}=(1-t^s)^{-1}.
\]
Hence its moments are $m!$ at multiples of $s$ and zero otherwise.
In a multinomial expansion every individual term is
$m!\zeta^{\sum_jjk_j}$, because the exponent factorials cancel the
multinomial denominator. Thus all terms have the same valuation
at every finite place of $\mathbb Q(\zeta)$, even in a nonzero moment.
An argument requiring a uniquely minimal term cannot prove SFC
universally. It must handle residual cancellation. This is the known
sharp linear example, not a counterexample or a new support theorem.

\section{What remains}
Neither these obstructions nor the existing critical-value argument
excludes a common nonzero projective zero of the actual $s$ moment
forms for every exponent matrix and starting index.
Strict Hankel or cumulant positivity is also insufficient, as shown in
\href{GAMMA_CONVOLUTION_APPELL_OBSTRUCTION.tex}{the retained
Gamma-convolution obstruction}; that example is not an ordinary
factorial phase with the required small support.

A proposed transfer of polynomial-composition dynamics must produce
an actual map with the required critical-value bound. A period rank
or a parameter-discriminant count alone is not that realization.
A possible Gamma-product Appell simplicity principle remains unproved,
and would still need a separate bridge to arbitrary polynomial supports.

The self-contained
\href{verify_unrestricted_sfc_mechanisms.py}{recurrence checker}
verifies 31 instances of (1), 30 of (2), growth controls and rejection
of a corrupted coefficient using exact integers. The all-index
identities, minimality and automorphism rigidity are the written proofs,
not consequences of those finite tests. Independent specialist review,
formal verification and a resolution of unrestricted SFC are not claimed.
\end{document}
